\ifdefined\OTCOMBINED%
\else
\documentclass{otscience}
% ============================================================
% Shared setup for OT Science modular workbook parts
% Place these files in the same folder as your fixed otscience.cls
% ============================================================

\makeatletter
\makeatother
\usepackage{multicol}
\usepackage{longtable}
\usepackage{makecell}
\usepackage{pdflscape}
\usepackage{bm}
\providecommand{\blank}[1][3cm]{\rule{#1}{0.4pt}}
\providecommand{\bigblank}{\rule{7cm}{0.4pt}}
\providecommand{\componentblank}{\blank[2.5cm]}
\providecommand{\R}{\mathbb{R}}
\providecommand{\C}{\mathbb{C}}
\providecommand{\ii}{\mathrm{i}}
\providecommand{\ee}{\mathrm{e}}
\providecommand{\dd}{\mathrm{d}}
\providecommand{\grad}{\nabla}
\providecommand{\divergence}{\nabla\!\cdot}
\providecommand{\curl}{\nabla\!\times}
\providecommand{\lap}{\nabla^2}
\providecommand{\ihat}{\hat{\mathbf{i}}}
\providecommand{\jhat}{\hat{\mathbf{j}}}
\providecommand{\khat}{\hat{\mathbf{k}}}
\providecommand{\er}{\hat{\mathbf{e}}_r}
\providecommand{\etheta}{\hat{\mathbf{e}}_\theta}
\providecommand{\ephi}{\hat{\mathbf{e}}_\phi}
\providecommand{\erho}{\hat{\boldsymbol{\rho}}}
\providecommand{\vA}{\mathbf{A}}
\providecommand{\vB}{\mathbf{B}}
\providecommand{\va}{\mathbf{a}}
\providecommand{\vb}{\mathbf{b}}
\providecommand{\vc}{\mathbf{c}}
\providecommand{\vx}{\mathbf{x}}
\providecommand{\vr}{\mathbf{r}}
\providecommand{\matA}{\mathbf{A}}
\providecommand{\matB}{\mathbf{B}}
\providecommand{\tr}{\operatorname{tr}}
\providecommand{\cof}{\operatorname{cof}}
\providecommand{\dev}{\operatorname{dev}}
\providecommand{\diag}{\operatorname{diag}}
\providecommand{\questionline}[1][5cm]{\rule{#1}{0.4pt}}
\setlength{\parindent}{0pt}
\setlength{\parskip}{4pt}
\renewcommand{\arraystretch}{1.25}

% Fallbacks in case the current otscience.cls has not defined nosplit boxes.
\makeatletter
\@ifundefined{formulaboxnosplit}{\newenvironment{formulaboxnosplit}[1][Formula]{\begin{formulabox}[#1]}{\end{formulabox}}}{}
\@ifundefined{definitionboxnosplit}{\newenvironment{definitionboxnosplit}[1][Definition]{\begin{definitionbox}[#1]}{\end{definitionbox}}}{}
\@ifundefined{summaryboxnosplit}{\newenvironment{summaryboxnosplit}[1][Summary]{\begin{summarybox}[#1]}{\end{summarybox}}}{}
\@ifundefined{warningboxnosplit}{\newenvironment{warningboxnosplit}[1][Warning]{\begin{warningbox}[#1]}{\end{warningbox}}}{}
\@ifundefined{exampleboxnosplit}{\newenvironment{exampleboxnosplit}[1][Example]{\begin{examplebox}[#1]}{\end{examplebox}}}{}
\makeatother

\newenvironment{practicebox}[1][Quick Drills and Practice]{\begin{examplebox}[#1]}{\end{examplebox}}
\newenvironment{practiceboxnosplit}[1][Quick Drills and Practice]{\begin{exampleboxnosplit}[#1]}{\end{exampleboxnosplit}}

\newcommand{\standalonetitle}[3]{%
    \title{\textbf{#1}\\\large #2}
    \author{OT Science Notes}
    \date{#3}
}

\standalonetitle{Symmetric and Skew-Symmetric Tensors, Invariants and Decompositions}{Part 5 of the OT Science Formula Completion Workbook}{\DTMtoday}
\begin{document}
\maketitle
\tableofcontents
\newpage
\fi

\section{Symm and Skew-Symm Tensors, Invariants  n' Decomps}

Second-order tensors appear naturally as stress tensors, strain tensors, inertia tensors, conductivity tensors, deformation gradients and vector gradients. Their symmetric and skew-symmetric parts often carry different physical meanings.

\subsection{Symmetric and Skew-Symmetric Tensors}

A tensor is symmetric if transposing it does not change it. A tensor is skew-symmetric if transposing it changes its sign. In three dimensions, a skew-symmetric tensor has only three independent components and can be represented by an axial vector.

\begin{formulaboxnosplit}[Symmetry Conditions]
    \[
        S_{ij}=S_{ji}\qquad\text{symmetry}.
    \]

    \[
        W_{ij}=-W_{ji}\qquad\text{skew-symmetry}.
    \]

    \[
        W_{ii}=0.
    \]

    \[
        A_{ij}=S_{ij}+W_{ij}.
    \]

    \[
        S_{ij}=\frac12(A_{ij}+A_{ji}),
        \qquad
        W_{ij}=\frac12(A_{ij}-A_{ji}).
    \]
\end{formulaboxnosplit}

\subsection{Connection with Vector Gradient}

If \(\mathbf{v}\) is a velocity field, then \(\nabla\mathbf{v}\) can be decomposed into symmetric and skew-symmetric parts. The symmetric part is related to strain rate, while the skew-symmetric part is related to local rotation.

\begin{formulabox}[Gradient Decomposition]
    \[
        L_{ij}=\partial_{j}v_i.
    \]

    \[
        D_{ij}=\frac12(L_{ij}+L_{ji})
        \qquad\text{symmetric rate-of-deformation tensor}.
    \]

    \[
        \Omega_{ij}=\frac12(L_{ij}-L_{ji})
        \qquad\text{skew spin tensor}.
    \]

    \[
        L_{ij}=D_{ij}+\Omega_{ij}.
    \]

    \[
        \omega_i={(\nabla\times\mathbf{v})}_i=\varepsilon_{ijk}\partial_{j}v_k.
    \]
\end{formulabox}

\subsection{Principal Invariants of a Second-Order Tensor}

Invariants do not change under a rotation of coordinates. They describe intrinsic properties of the tensor.

\begin{formulaboxnosplit}[Principal Invariants]
    For a \(3\times3\) tensor \(\mathbf{A}\),
    \[
        I_1=\tr(A)=A_{ii}.
    \]

    \[
        I_2=\frac12\left[{(\tr A)}^2-\tr(A^2)\right].
    \]

    \[
        I_3=\det(A).
    \]

    The characteristic polynomial may be written as
    \[
        \det(\lambda I-A)=\lambda^3-I_1\lambda^2+I_2\lambda-I_3.
    \]
    Equivalently,
    \[
        \det(A-\lambda I)=-\lambda^3+I_1\lambda^2-I_2\lambda+I_3.
    \]
\end{formulaboxnosplit}

\subsection{Deviatoric and Spherical Parts}

Many physical tensors can be decomposed into a mean or spherical part and a deviatoric part. In stress analysis, the spherical part is associated with hydrostatic pressure, while the deviatoric part is associated with shape distortion.

\begin{formulabox}[Deviatoric Decomposition]
    \[
        \operatorname{sph}{(A)}_{ij}=\frac13A_{kk}\delta_{ij}.
    \]

    \[
        A'_{ij}=\dev{(A)}_{ij}=A_{ij}-\frac13A_{kk}\delta_{ij}.
    \]

    \[
        A_{ij}=\frac13A_{kk}\delta_{ij}+A'_{ij}.
    \]

    \[
        A'_{kk}=0.
    \]

    For a stress tensor \(\sigma_{ij}\),
    \[
        p=-\frac13\sigma_{kk},
        \qquad
        s_{ij}=\sigma_{ij}-\frac13\sigma_{kk}\delta_{ij}.
    \]
\end{formulabox}

\subsection{Invariants of Deviatoric Tensors}

For a deviatoric tensor \(A'\), the first invariant is zero. The remaining invariants are important in mechanics.

\begin{formulabox}[Deviatoric Invariants]
    \[
        J_1=\tr(A')=0.
    \]

    \[
        J_2=\frac12A'_{ij}A'_{ij}.
    \]

    \[
        J_3=\det(A').
    \]

    For stress deviator \(s_{ij}\), one common convention is
    \[
        J_2=\frac12s_{ij}s_{ij}.
    \]
\end{formulabox}

\begin{exampleboxnosplit}[Worked Example: Tensor Decomposition]
    Let
    \[
        A=\begin{pmatrix}
            2 & 1 & 3  \\
            4 & 0 & -1 \\
            5 & 2 & 1
        \end{pmatrix}.
    \]
    Then
    \[
        S=\frac12(A+A^T),
        \qquad
        W=\frac12(A-A^T).
    \]
    Compute
    \[
        S=\begin{pmatrix}
            2   & 5/2 & 4   \\
            5/2 & 0   & 1/2 \\
            4   & 1/2 & 1
        \end{pmatrix},
        \qquad
        W=\begin{pmatrix}
            0   & -3/2 & -1   \\
            3/2 & 0    & -3/2 \\
            1   & 3/2  & 0
        \end{pmatrix}.
    \]
    Check that \(S^T=S\), \(W^T=-W\), and \(A=S+W\).
\end{exampleboxnosplit}

\begin{practicebox}[Quick Drills and Practice: Tensor Decompositions and Invariants]
    \textbf{Level A:\,Symmetry checks}
    \begin{multicols}{2}
        \begin{enumerate}
            \item Is \(\begin{pmatrix}1&2\\2&3\end{pmatrix}\) symmetric?
            \item Is \(\begin{pmatrix}0&-4\\4&0\end{pmatrix}\) skew-symmetric?
            \item Complete: \(S_{ij}-S_{ji}=\blank[2cm]\).
            \item Complete: \(W_{ij}+W_{ji}=\blank[2cm]\).
            \item Complete: \(W_{ii}=\blank[2cm]\).
            \item Complete: \(A_{ij}=\blank[3cm]\).
            \item Complete: \(S_{ij}=\blank[3cm]\).
            \item Complete: \(W_{ij}=\blank[3cm]\).
        \end{enumerate}
    \end{multicols}

    \textbf{Level B:\,Decomposition}
    \begin{enumerate}
        \item Decompose \(A=\begin{pmatrix}1&2&3\\0&4&5\\1&0&6\end{pmatrix}\) into symmetric and skew-symmetric parts.
        \item Decompose \(A=\begin{pmatrix}0&-1&2\\1&0&-3\\-2&3&0\end{pmatrix}\). What do you notice?
        \item For \(\mathbf{v}=(-y,x,0)\), compute \(\nabla\mathbf{v}\), its symmetric part and its skew part.
        \item For \(\mathbf{v}=(ax,by,cz)\), compute \(\nabla\mathbf{v}\), its trace and its symmetric/skew parts.
    \end{enumerate}

    \textbf{Level C:\,Invariants}
    \begin{enumerate}
        \item Compute \(I_1,I_2,I_3\) for \(A=\diag(1,2,3)\).
        \item Compute \(I_1,I_2,I_3\) for \(A=\diag(a,b,c)\).
        \item Compute \(I_1,I_2,I_3\) for \(A=\begin{pmatrix}2&1&0\\1&2&0\\0&0&3\end{pmatrix}\).
        \item Write the characteristic polynomial of a \(3\times3\) tensor using \(I_1,I_2,I_3\).
        \item If a tensor has eigenvalues \(\lambda_1,\lambda_2,\lambda_3\), express \(I_1,I_2,I_3\) in terms of the eigenvalues.
    \end{enumerate}

    \textbf{Level D:\,Deviatoric tensors}
    \begin{enumerate}
        \item Find the deviatoric part of \(A=\diag(4,1,1)\).
        \item Find the deviatoric part of \(A=\diag(6,3,0)\).
        \item Prove that \(A'_{kk}=0\).
        \item For \(A=\diag(4,1,1)\), compute \(J_2=\frac12A'_{ij}A'_{ij}\).
        \item For \(\sigma=\diag(100,50,25)\), compute the hydrostatic part and deviatoric part.
    \end{enumerate}

    \textbf{Level E:\,Challenge}
    \begin{enumerate}
        \item Show that the trace of a skew-symmetric tensor is zero.
        \item Show that \(x_{i}W_{ij}x_j=0\) for any skew-symmetric tensor \(W_{ij}\).
        \item Show that \(x_{i}S_{ij}x_j\) depends only on the symmetric part of a tensor.
        \item If \(A\) is symmetric, explain why its eigenvalues are real in ordinary Euclidean space.
        \item Explain why invariants are useful when comparing the same tensor in two different coordinate frames.
    \end{enumerate}
\end{practicebox}

\ifdefined\OTCOMBINED%
\else
\end{document}
\fi
